\chapter*{ПРИЛОЖЕНИЕ А}
\addcontentsline{toc}{chapter}{ПРИЛОЖЕНИЕ А. Реестр экспериментальных материалов}
\begin{center}\bfseries Реестр экспериментальных материалов\end{center}
\noindent Программный стенд и журналы находятся в \texttt{code/}; документация и литература --- в \texttt{materials/}; результаты серий и материалы отчёта --- в \texttt{nir/}; материалы будущей ВКР --- в \texttt{vkr/}.

\begin{itemize}
\item протоколы и source receipts: \texttt{code/runs/};
\item исходные таблицы, графики и отчёты: \texttt{nir/results/};
\item реестр источников и локальные полные тексты: \texttt{materials/literature/};
\item карта переноса и проверка сохранности: \texttt{materials/repository\_layout/}.
\end{itemize}

Точные значения, определения эффектов, строки таблиц и контрольные суммы следует проверять по исходным артефактам~\cite{v4tests,v4audit}. Начальное наполнение настоящего шаблона не заменяет полный итоговый текст отчёта НИР.
